\section{Effective resolutions and fiberwise nonvanishing}\label{sec:resolution-fibers}

Fujiki class~$\mathcal C$ consists of compact complex spaces
bimeromorphic to compact Kähler manifolds. For a rational adjoint,
$\kappa$ denotes its Kodaira dimension, defined by sections of divisible
positive multiples and equal to $-\infty$ when all such section spaces
vanish. A fibration is a surjective holomorphic map with connected
fibres.

We first make the boundary effective on a resolution while preserving
the divisible adjoint section spaces. We then establish fiberwise
nonvanishing with one common divisible degree.

\begin{lemma}\label{lem:supply-log-resolution}
Let $(X,B)$ be a rational klt pair with $D=K_X+B$ rational Cartier, and let
$p:Z\to X$ be a projective log resolution.  There is an effective rational
SNC divisor $\Gamma$ with coefficients strictly less than one such that
\begin{equation}\label{eq:supply-log-resolution}
 K_Z+\Gamma\simq p^*D+E,
 \qquad E\geq0\quad\text{and}\quad E\text{ is }p\text{-exceptional}.
\end{equation}
These are actual rational line bundles with their birational
identification.  If $D$ is analytically pseudo-effective, then so is
$K_Z+\Gamma$.  For every sufficiently divisible positive integer $m$,
\begin{equation}\label{eq:supply-section-descent}
 p_*\cO_Z\bigl(m(K_Z+\Gamma)\bigr)
 \simeq \cO_X(mD).
\end{equation}
\end{lemma}

\begin{proof}
Define the crepant subboundary $\Gamma_0$ by
$K_Z+\Gamma_0=p^*D$, using the canonical identification over the locus
where $p$ is an isomorphism.  Local canonical generators and their
Jacobians define this equality without choosing a global canonical
divisor.  Klt singularities give
$\operatorname{coeff}_P\Gamma_0<1$ for every prime divisor $P$ on $Z$.
The nonexceptional coefficients are those of $B$ and are nonnegative.
Set
\[
 \Gamma=\sum_P\max\{\operatorname{coeff}_P\Gamma_0,0\}P,
 \qquad E=\Gamma-\Gamma_0.
\]
Both supports lie in the SNC divisor of the log resolution.  This proves
\eqref{eq:supply-log-resolution}.  A positive singular metric on a
Cartier multiple of $D$ pulls back under the dominant map $p$; multiplying
by the divisor metric of the effective correction proves
pseudo-effectivity upstairs.

Take $m$ clearing the Cartier indices and all coefficients in
\eqref{eq:supply-log-resolution}.  The image of $\Supp E$ has codimension
at least two in the normal space $X$, and
$p_*\cO_Z(mE)=\cO_X$.  Indeed, a local section of $\cO_Z(mE)$ is a
meromorphic function with poles allowed only along $E$.  It gives a
holomorphic function off that codimension-two image, which extends by
normality; its pullback agrees with the original section on a dense open
set.  The reverse inclusion is immediate.  The actual line-bundle
identity
\[
 \cO_Z\bigl(m(K_Z+\Gamma)\bigr)
 \simeq p^*\cO_X(mD)\otimes\cO_Z(mE)
\]
and the projection formula give \eqref{eq:supply-section-descent}.
\end{proof}

The next lemma separates two quantifiers in restricting a positive
metric.  Such a metric initially restricts on almost every smooth
fiber.  Projective nonvanishing in dimension at most three and proper
holomorphic cohomology then supply a single nonzero adjoint multiple
on every fiber of the indicated smooth locus.

\begin{lemma}\label{lem:supply-fibers}
Let $g:Z\to T$ be a surjective holomorphic map with connected fibers
between smooth compact connected Kähler manifolds.  Let $\Gamma$ be an
effective rational SNC divisor with coefficients strictly less than one.
Suppose that $K_Z+\Gamma$ is analytically pseudo-effective and that the
very general smooth fiber is projective of dimension at most three.
Then
\[
 \kappa(F,K_F+\Gamma|_F)\geq0
\]
for a very general smooth fiber $F$.  Moreover, one positive divisible
degree gives a nonzero section on every fiber over the smooth base locus
where the boundary strata restrict smoothly.
\end{lemma}

\begin{proof}
\emph{Restriction of the metric.}
Remove the critical values of $g$ and the proper images on which a
boundary stratum fails to restrict smoothly.  The remaining nonempty
open set $T^\circ\subset T$ is connected.  On it the fibers are smooth,
$\Gamma|_{F_t}$ is an effective SNC boundary with coefficients less than
one, and adjunction is an identity of actual rational line bundles:
\begin{equation}\label{eq:supply-fiber-adjunction}
 (K_Z+\Gamma)|_{F_t}=K_{F_t}+\Gamma|_{F_t}.
\end{equation}
Here the constant one-dimensional factor coming from
$K_T|_t$ is immaterial to the line-bundle isomorphism on $F_t$.

Choose $r>0$ so that $L=\cO_Z(r(K_Z+\Gamma))$ is a line bundle.  Give $L$
a singular Hermitian metric of nonnegative curvature.  Its local weights
are plurisubharmonic.  In product coordinates for the submersion over
$T^\circ$, their local integrability and Fubini's theorem show that the
weights restrict to plurisubharmonic functions, rather than identically
$-\infty$, on almost every fiber.  A countable cover gives this assertion
simultaneously for the local weights; their transition identities are
preserved under restriction.  Thus $L|_{F_t}$ is analytically
pseudo-effective for almost every $t\in T^\circ$.

\smallskip\noindent
\emph{Projective nonvanishing on the fibers.}
For almost every such parameter the fiber is also projective: the
exceptions to the very-general projectivity assumption have measure
zero.  On a smooth projective manifold the analytic pseudo-effective
cone restricted to divisor classes is the closure of the effective
divisor cone, by \cite[Proposition~1.2]{BDPP2013}.  Hence the projective
klt pair $(F_t,\Gamma|_{F_t})$ has pseudo-effective adjoint.  The
projective log MMP in dimensions at most three produces a nef model.
For example, in dimension three the ordinary cone and flip theorems,
together with \cite[Theorem~1.2]{CT2023} with zero nef part, give
termination; pseudo-effectivity excludes a Mori fiber space.  The
projective log abundance theorem
\cite[Theorem~1.1]{KMMAbundance1994}, with the correction
\cite[\S6.1]{MatsukiCorrection2003}, gives a nonzero section of a
multiple of the nef adjoint.  Dimensions one and two use the classical
log abundance theorem.  The effective pullback comparison on a common
resolution pulls this section back to a section of a multiple of the
original adjoint.  It descends to $F_t$ by the projection formula for the
modification.  Consequently
\begin{equation}\label{eq:supply-ae-nonvanishing}
 H^0(F_t,L^{\otimes m_t}|_{F_t})\ne0
 \quad\text{for some }m_t>0
\end{equation}
for almost every $t\in T^\circ$.  If the fiber is a point this assertion
holds directly.

\smallskip\noindent
\emph{A common positive degree.}
To obtain the required quantifier, put
\[
 A_m=\bigl\{t\in T^\circ:
    h^0(F_t,L^{\otimes m}|_{F_t})\geq1\bigr\},
 \qquad m=1,2,\ldots .
\]
These are closed analytic jumping loci: the restricted family is proper
and smooth, and $L^{\otimes m}$ is locally free and flat over its base,
so proper holomorphic cohomology and base change apply
\cite[\S8, Theorem~3, p.~62]{Douady1974}.  By
\eqref{eq:supply-ae-nonvanishing}, their countable union contains almost
every point of $T^\circ$.  If every $A_m$ were proper, that union would
have measure zero.  Therefore some $A_m$ equals $T^\circ$, since this
base is connected.  This proves both assertions.  In particular,
almost-everywhere restriction of a metric has not been identified with
the very-general assertion by definition.
\end{proof}
